\subsection{Hermitian geometry.}

DEFINITION 2. --- Let A be a field, L an affine space over A, and T the space of translations of L (chap.~II, 2 \(^{e}\) ed., App. II). If T is equipped with a nondegenerate Hermitian form Φ, then L is said to be a Hermitian space over A, and Φ is called the metric form of L.

If J is the identity (which implies that A is commutative), one says instead that L is a Euclidean space.

If \(a\) and \(b\) are two points of L, set \(e(a, b) = \Phi(b - a, b - a)\) . Let \(c\) be a third point of L. For \(b - a\) and \(c - a\) to be orthogonal, it is necessary, according to formula (17) of § 1, no. 5, that one have \(e(b, c) = e(a, b) + e(a, c)\) and this condition is sufficient when J = 1 and A is not of characteristic 2 (“Pythagorean theorem”).

Two linear varieties of L are said to be orthogonal if their directions (chap.~II, 2 \(^{e}\) ed., App. II, no. 3) are orthogonal. A linear manifold of L is said to be isotropic (resp. totally isotropic) if its direction is isotropic (resp. totally isotropic). A vector of T is said to be orthogonal to a linear manifold of L if it is orthogonal to the direction of this manifold.

Let V be a linear variety in L, and let x be a point of L. The set of points y of L such that y − x is orthogonal to V is a linear variety W passing through x; we say that W is the totally orthogonal variety (or, more simply, orthogonal) to V passing through \(x\) . If L is finite-dimensional, the dimension of W is equal to the codimension of V. Moreover, if V is non-isotropic, the directions of V and W are supplementary (§ 4, no. 1, cor. of prop. 1); then W meets V in a single point \(x_1\) ; taking an origin in V, one sees immediately that, for fixed V, the map \(x \to x_1\) is an idempotent affine linear map; it is called the orthogonal projection of L onto V; the linear map associated with it (chap.~II, 2 \(^{e}\) ed., App. II, no. 4) is the orthogonal projector of T onto the direction of V (no. 3).

DEFINITION 3. --- Let L be a Hermitian space over a field A, and T the space of translations of L. One calls a displacement (resp. similarity) of L any affine bijection u of L onto L such that the linear map v associated with u in T (chap.~II, 2 \(^{e}\) ed., App. II, no. 4) is unitary (resp. is a similarity).

The group of translations is a distinguished subgroup of the affine group; it is therefore a distinguished subgroup of the group of similarities and of the group of displacements. For every \(a \in L\) , let \(G_a\) be the group of similarities (resp. displacements) leaving \(a\) fixed; if one identifies L with T by taking \(a\) as origin, \(G_a\) is the group of similarities (resp. the unitary group) of T. Every similarity (resp. displacement) \(u\) can be written, in one and only one way, in the form \(u = u_1 t_1\) where \(u_1 \in G_a\) and \(t_1 \in T\) , and also in the form \(u = t_2 u_2\) where \(u_2 \in G_a\) and \(t_2 \in T\) ; moreover, one has \(u_2 = u_1\) and \(t_2 = u_1 t_1 u_1^{-1}\) (chap.~II, 2 \(^\text{e}\) ed., App. II, no. 4).

Let \(u\) be a similarity in L, \(\varrho\) the associated similarity in T. The multiplier of \(\varrho\) is also called the multiplier of \(u\) (no. 5). If we denote by \(\alpha(u)\) this multiplier, the map \(u \to \alpha(u)\) is a homomorphism from the group of similarities of L into the multiplicative group of invertible elements of the center of A; its kernel is the group of displacements, which is therefore a normal subgroup of the group of similarities. When A is commutative and L is finite-dimensional, there is, between the determinant det \(u\) (equal by definition to det \(\varrho\) ) and \(\alpha(u)\) , the same relations as in no. 5. The displacements \(u\) such that det \(u = 1\) form a distinguished subgroup of the group of displacements; this subgroup has index 2 if A is a commutative field of characteristic ≠ 2 and J is the identity.

PROPOSITION 6. --- Let L be a finite-dimensional Hermitian space over A, whose metric form has index 0. Every similarity u of L, with multiplier \(\mu \neq 1\) then admits one and only one fixed point.

Indeed, let \(a\) be a point of L. There exists a similarity \(\nu\) of L leaving \(a\) fixed and a translation \(t\) of L such that \(u = tv\) . To say that \(b\) is a fixed point of \(u\) amounts to saying that \(\nu(b) - b = t\) . To show that this equation admits one and only one solution \(b\) , let us identify L with its space of translations T by taking \(a\) as origin. It then suffices to prove that the endomorphism \(\nu - 1\) of T is invertible, in other words that the relation \(\nu(x) - x = 0\) ( \(x \in \mathrm{T}\) ) implies \(x = 0\) . Now, if \(\nu(x) - x = 0\) , one has \(\Phi(x, x) = \Phi(\nu(x), \nu(x)) = \mu \Phi(x, x)\) , hence \(\Phi(x, x) = 0\) since \(\mu \neq 1\) ; this entails \(x = 0\) since \(\Phi\) has index 0. QED.

Suppose that A is a field of characteristic \(\neq 2\) Every displacement \(u\) of L such that \(u^2 = 1\) admits at least one fixed point, for example the midpoint \(\frac{1}{2}(x + u(x))\) of two homologous points; taking this point as origin, one sees that the unitary automorphism of T associated with \(u\) is a symmetry (no. 3). Let V be a non-isotropic linear variety in L; one says that a displacement \(u\) is the symmetry with respect to V if, taking an origin in V, \(u\) is identified with the symmetry of T with respect to V. It is equivalent to say that \(u(x)\) is obtained as follows: denoting by \(x_1\) the orthogonal projection of \(x\) on V, we have \(u(x) - x = 2(x_1 - x)\) .

Exercises. --- 1) Suppose that A is a commutative field. Given a Hermitian matrix R of order n over A, the principal minors of order r of R are called the minors obtained by deleting from R n - r rows and the n - r columns with the same indices.

\begin{enumerate}
\def\labelenumi{\alph{enumi})}
\setcounter{enumi}{1}
\tightlist
\item
  From a), deduce that if R has rank r, there exists a permutation \(\sigma \in S_{n}\) such that, if one applies the same permutation \(\sigma\) to the rows and columns of R and if we denote by S the matrix thus obtained, by \(\Delta_{k}\) the principal minor of order k of S obtained by deleting from S the rows and columns with index \textgreater{} k , one has the following two properties: \(1^{\circ} \Delta_{r} \neq 0\) ; \(2^{\circ}\) there is no index k \textless{} r such that \(\Delta_{k} = \Delta_{k+1} = 0\) .
\end{enumerate}

\begin{enumerate}
\def\labelenumi{\arabic{enumi})}
\setcounter{enumi}{1}
\tightlist
\item
  We assume that A is a commutative field, and that E is of finite dimension n. Let \(\Phi\) be a Hermitian sesquilinear form on E, satisfying condition (T) of § 4, no. 2, \(R = (\alpha_{ij})\) the matrix of \(\Phi\) with respect to a basis ( \(e_i\) ) of E.
\end{enumerate}

\begin{enumerate}
\def\labelenumi{\alph{enumi})}
\item
  If \(\Phi\) has rank \(r\) and if the principal minor (exerc. 1) obtained by deleting in \(R\) the rows and columns with indices \(>r\) is not zero, show that there exists a new basis \((f_i)\) of E such that \(e_i = f_i\) for \(1 \leqslant i \leqslant r\) and that the matrix of \(\Phi\) with respect to \((f_i)\) is obtained by replacing by 0 in \(R\) all the \(\alpha_{ij}\) such that \(i > r\) or \(j > r\) (consider the subspace E \(^\circ\) orthogonal to E).
\item
  From a), deduce that if \(\Phi\) has rank \(n\) , and if the cofactor \(\Delta_{n-1}\) of \(\alpha_{nn}\) in the determinant \(\Delta = \det R\) is nonzero, there exists a new basis \((f_i)\) of E such that \(f_i = e_i\) for \(1 \leqslant i \leqslant n-1\) , and such that one has
\end{enumerate}

\[
\Phi (x, y) = \Phi \left(\sum_ {i = 1} ^ {n} \xi_ {i} f _ {i}, \sum_ {i = 1} ^ {n} \eta_ {i} f _ {i}\right) = \sum_ {i = 1} ^ {n - 1} \sum_ {j = 1} ^ {n - 1} \alpha_ {i j} \xi_ {i} \bar {\eta} _ {j} + \frac {\Delta}{\Delta_ {n - 1}} \xi_ {n} \bar {\eta} _ {n}
\]

(consider the Hermitian form whose matrix with respect to \((e_i)\) is obtained by replacing \(\alpha_{nn}\) by \(\alpha_{nn} - \frac{\Delta}{\Delta_{n-1}}\) in \(R\) ).

\begin{enumerate}
\def\labelenumi{\alph{enumi})}
\setcounter{enumi}{2}
\tightlist
\item
  We assume that \(\Phi\) has rank \(n\) , that \(\Delta_{n-1} = 0\) , but the principal minor \(\Delta_{n-2}\) of \(R\) obtained by deleting the rows and columns with indices \(n-1\) and \(n\) in \(R\) is nonzero. Show that there exists a new basis \((f_i)\) of E such that \(f_i = e_i\) for \(1 \leqslant i \leqslant n-2\) , and such that one has
\end{enumerate}

\[
\Phi (x, y) = (\sum_ {i = 1} ^ {n} \xi_ {i} f _ {i}, \sum_ {i = 1} ^ {n} \eta_ {i} f _ {i}) = \sum_ {i = 1} ^ {n - 2} \sum_ {j = 1} ^ {n - 2} \alpha_ {i j} \xi_ {i} \bar {\eta} _ {j} + \xi_ {n - 1} \bar {\eta} _ {n} + \xi_ {n} \bar {\eta} _ {n - 1}.
\]

(If H is the hyperplane spanned by \(e_{1},\ldots,e_{n-1}\) which is isotropic, note that the line orthogonal to H is not contained in the subspace generated by \(e_{1},\ldots,e_{n-2}\) , and use prop. 2 of § 4, no. 2).

\begin{enumerate}
\def\labelenumi{\arabic{enumi})}
\setcounter{enumi}{2}
\item
  Let A be a finite field, E a finite-dimensional vector space over A, Φ a nondegenerate Hermitian sesquilinear form on E, relative to an automorphism J ≠ 1 of A. Show that E admits an orthonormal basis for Φ (cf.~chap.~V, § 11, no. 5, cor. of th. 3).
\item
  Let A be a finite field of characteristic \(\neq 2\) , E a finite-dimensional vector space of dimension \(n\) over A.
\end{enumerate}

\begin{enumerate}
\def\labelenumi{\alph{enumi})}
\item
  Show that for every nondegenerate symmetric bilinear form \(\Phi\) on E, there exists an orthogonal basis \((e_i)\) of E such that \(\Phi(e_i, e_i) = 1\) for \(1 \leqslant i \leqslant n - 1\) , \(\Phi(e_n, e_n) = \Delta\) (discriminant of \(\Phi\) with respect to \((e_i)\) ). (Note that if \(\alpha\beta \neq 0\) , the equation \(\alpha\xi^2 + \beta\eta^2 = \gamma\) always admits solutions \((\xi, \eta)\) in A if \(\gamma \neq 0\) (chap.~V, § 11, exerc. 4)).
\item
  For two nondegenerate symmetric bilinear forms on E to be equivalent, it is necessary and sufficient that the ratio of their discriminants (with respect to the same basis of E) be a square in A. Deduce that, if n is odd, for every symmetric bilinear form \(\Phi\) nondegenerate on E, there exists an orthogonal basis with respect to which the matrix of \(\Phi\) is of the form \(\lambda I_{n} (\lambda \in \mathbf{A})\) ; the index of \(\Phi\) is then \((n - 1)/2\) .
\item
  If \(n=2m\) is even, show that the index of a nondegenerate symmetric bilinear form \(\Phi\) on E is \(m\) if \((-1)^{m}\Delta\) is a square in A, \(m-1\) otherwise.
\end{enumerate}

\begin{enumerate}
\def\labelenumi{\arabic{enumi})}
\setcounter{enumi}{4}
\tightlist
\item
  Let A be a commutative field of characteristic \(\neq 2\) Let I be a polynomial with coefficients in A, in the \(n(n+1)/2\) indeterminates \(X_{ij}\) ( \(1 \leqslant i \leqslant j \leqslant n\) ); for every symmetric matrix \(R = (\alpha_{ij})\) over a commutative extension field A' of A, denote by I(R) the element of A' obtained by substituting \(\alpha_{ij}\) to the indeterminate \(X_{ij}\) ( \(i \leqslant j\) ) in I.
\end{enumerate}

We assume that I is such that, for the matrix \(U = (u_{ij})\) with \(u_{ij} = \mathrm{X}_{ij}\) for \(i \leqslant j\) , \(u_{ij} = \mathrm{X}_{ji}\) for \(i > j\) and the square matrix \(P = (\mathrm{Y}_{ij})\) of order \(n\) (where the \(\mathrm{Y}_{ij}\) are \(n^2\) other indeterminates), one has

\[
\mathrm{I} (\mathbf {t} P U P) = (\det P) ^ {\hbar} \mathrm{I} (U)
\]

where h is an integer \textgreater{} 0. Show that h is even and that I(U) = γ(det U)k, where h = 2k and γ ∈ A. (Using th. 1, show that for every symmetric matrix R over the algebraic closure Ω of A, one has (I(R))² = λ(det R)h, where λ ∈ Ω, and use the fact that the polynomial det U with respect to the Xij is not a square, by considering the terms of this polynomial containing an Xii).

*6) Let A be a complete non-discrete valued field, commutative and of characteristic \(\neq 2\) (Top. gen., chap.~IX, § 3, no. 2), \(\Phi\) a non-degenerate Hermitian form on a finite-dimensional vector space E of dimension \(n\) over A, \(R = (\alpha_{ij})\) the matrix of \(\Phi\) with respect to a basis \((e_i)\) of E. Show that there exists \(\varepsilon > 0\) such that, for every Hermitian matrix \(R' = (\alpha_{ij}')\) satisfying the conditions \(|\alpha_{ij}' - \alpha_{ij}| \leqslant \varepsilon\) for every pair \((i, j)\) , the form \(\Phi'\) having \(R'\) as matrix with respect to the basis \((e_i)\) is equivalent to \(\Phi\) . (Reduce to the case where \(R\) is diagonal; use exerc. 2 b) relying on the following lemma: there exists a number \(a > 0\) such that for \(|\eta| \leqslant a\) , there exists in A an element \(\xi\) such that \(\xi^2 = 1 - \eta\) . To prove this lemma, use the binomial series for \((1 - x)^{1/2})\cdot *\)

¶ 7) Let A be a non-orderable commutative field (chap.~VI, § 2, exerc. 8) of characteristic ≠ 2, E a vector space of finite dimension n \textgreater{} 0 over A, Q a non-degenerate quadratic form on E, \((e_i)\) an orthogonal basis for Q, so that \(\mathrm{Q}\left(\sum_{i=1}^{n}\xi_{i}e_{i}\right)=\sum_{i=1}^{n}\alpha_{i}\xi_{i}^{2}\) . For \(1\leqslant r\leqslant n\) , set \(\mathrm{Q}_{r}(\xi_{1},\ldots,\xi_{r})=\sum_{i=1}^{r}\alpha_{i}\xi_{i}^{2}\) , and denote by \(M_{r}\) the set of values of \(Q_{r}\) as the \(\xi_{i}\left(1\leqslant i\leqslant r\right)\) range over A.

\begin{enumerate}
\def\labelenumi{\alph{enumi})}
\item
  Show that if, for some index \(r\) , one has \(\mathbf{M}_r = \mathbf{M}_{r+1}\) , then it follows that \(\mathbf{M}_r = \mathbf{A}\) (note that every element of A is a sum of squares (chap.~VI, § 2, exerc. 7)).
\item
  Suppose that the subgroup S of the multiplicative group A*, formed by the squares of elements of A, has finite index s in A*. Deduce from a) that if n \textgreater{} s, every non-degenerate quadratic form on E has index \textgreater{} 0 (note that every set M \(_{r}\) is the union of 0 and of classes mod. S). *Application to the case where A is a p-adic field Q \(_{p}\) (Top. gen., chap.~III, § 5, exerc. 35).*
\end{enumerate}

\begin{enumerate}
\def\labelenumi{\arabic{enumi})}
\setcounter{enumi}{7}
\item
  Let A be a commutative field of characteristic \(\neq 2\) , E a finite-dimensional vector space \(n\) over A, Q a non-degenerate quadratic form of index 0 on E. Let A' be an algebraic extension of A, of finite odd degree, E' the vector space over A' obtained by extending to A' the field of scalars of E. Show that the extension Q' of Q to E' (§ 3, no. 4, prop. 3) still has index 0. (Reduce to the case where A' = A{[}X{]}/(f), f being an irreducible polynomial of odd degree m over A. Let ( \(e_i\) ) be an orthogonal basis of E for Q, and \(\rho_i = Q(e_i)\) ; show that, in A{[}X{]}, a relation of the form \(\sum_{i} \rho_i(g_i(X))^2 = f(X)h(X)\) , where the \(g_i\) are polynomials, not all zero, of degree \(\leqslant m - 1\) , is impossible; for this observe that h would necessarily have odd degree, and consider an irreducible factor of h, of odd degree).
\item
  Let A be a field, E a vector space over A admitting a countable basis \((e_n)_{n\geqslant 1}\) , \(\Phi\) a non-degenerate Hermitian sesquilinear form on E, satisfying condition (T) ( \(\S 4\) , no. 2).
\end{enumerate}

\begin{enumerate}
\def\labelenumi{\alph{enumi})}
\item
  Show that if the conditions (C) of th. 1 are not simultaneously satisfied, there exists in E an orthogonal basis for \(\Phi\) (reason as in exerc. 4 of § 5).
\item
  Suppose moreover that A is commutative, and that there exists an integer s such that on every finite-dimensional vector space of dimension \textgreater s over A, every non-degenerate Hermitian sesquilinear form has index \textgreater0 (cf.~exerc. 7). Show that there then exists in E an orthonormal basis for Φ. (Reason as in a), observing that for every element of A of the form α = λ + λ̄, and every non-degenerate Hermitian form Ψ on a finite-dimensional space F of dimension \textgreater s, there exists z ∈ F such that Ψ(z, z) = α (cf.~§ 4, no. 2, prop. 4).)
\end{enumerate}

¶ 10) a) Let A be a principal ideal ring in which there is only one maximal ideal Aπ, such that 2 is not divisible by π (chap.~VII, § 1, exerc. 4). Let E be a free module over A, of dimension n.~Show that every symmetric bilinear form Φ on E admits an orthogonal basis. (Let r be the largest exponent such that πr divides all elements Φ(x, y); show that there exists a ∈ E such that Φ(a, a) = απr, where α is invertible in A; deduce that E is the direct sum of F = Aa and of the submodule F⁰ orthogonal to F.)

\begin{enumerate}
\def\labelenumi{\alph{enumi})}
\setcounter{enumi}{1}
\item
  Give an example (for n = 2) where \(\Phi\) is non-degenerate and where there exists a non-isotropic submodule F of E, of rank 1, admitting a complement in E but such that \(F^{0}\) is not a complement of F.
\item
  Let \((e_i)\) be an orthogonal basis for \(\Phi\) , and \(\alpha_i = \Phi(e_i, e_i)\) . Show that the ideals \(\mathbf{A}\alpha_i\) are, up to order, independent of the orthogonal basis considered (cf.~§ 5, th. 1).
\end{enumerate}

These ideals are said to be the invariant factors of the form \(\Phi\) . Give an example of two forms having the same invariant factors and not equivalent (take two forms whose quotient of discriminants is not a square).

\begin{enumerate}
\def\labelenumi{\alph{enumi})}
\setcounter{enumi}{3}
\item
  Let F be a submodule of E, \(\Phi_{\mathrm{F}}\) the restriction of \(\Phi\) to \(\mathbf{F} \times \mathbf{F}\) , \(\mathbf{A}\alpha_{i}\) ( \(1 \leqslant i \leqslant r\) ) the nonzero invariant factors of \(\Phi\) , arranged so that \(\alpha_{i}\) divides \(\alpha_{i+1}\) , \(\mathrm{A}\beta_{i}\) ( \(1 \leqslant i \leqslant s\) ) the nonzero invariant factors of \(\Phi_{\mathrm{F}}\) , arranged so that \(\beta_{i}\) divides \(\beta_{i+1}\) . Show that we have \(s \leqslant r\) and \(\beta_{i}\) is a multiple of \(\alpha_{i}\) for \(1 \leqslant i \leqslant s\) (same method as in exerc. 1 a) of § 5).
\item
  We assume \(\Phi\) nondegenerate; let F, G be two non-isotropic submodules of E such that \(F^{0}\) (resp. \(G^{0}\) ) is a supplement of F (resp. G). One assumes that the restrictions of \(\Phi\) to F and to G are equivalent; show that there then exists an automorphism \(u\) of E, leaving invariant \(\Phi\) , and such that \(u(F) = G\) . (Using \(a\) ), one can reduce to the case where \(F = Aa\) , \(G = Ab\) , \(\Phi(a, a) = \Phi(b, b)\) . Let \((c_{j})\) be a basis of \(G^{0}\) , and let \(b'\) , \(c_{j}'\) ( \(1 \leqslant j \leqslant n - 1\) ) be the components of \(b\) and \(c_{j}\) respectively in \(F^{0}\) ; show that there exist scalars \(\mu_{j}\) ( \(1 \leqslant j \leqslant n - 1\) ) such that the elements \(d_{j} = c_{j}' + \mu_{j}b'\) satisfy the relations \(\Phi(d_{j}, d_{k}) = \Phi(c_{j}, c_{k})\) for every pair of indices; to do this, note that for every \(\lambda \in A\) one of the elements \(1 \pm \lambda\) is invertible in A.)
\end{enumerate}

\begin{enumerate}
\def\labelenumi{\arabic{enumi})}
\setcounter{enumi}{10}
\item
  Let A be a principal ideal ring of characteristic 0, in which there is only one principal ideal \(\pi\) , such that 2 is divisible by \(\pi\) . If \((e_1, e_2)\) is the standard basis of \(E = A^2\) , \(\Phi\) the symmetric bilinear form on E defined by \(\Phi(\xi_1 e_1 + \xi_2 e_2, \eta_1 e_1 + \eta_2 e_2) = \xi_1 \eta_2 + \xi_2 \eta_1\) , show that there does not exist an orthogonal basis of E for \(\Phi\) .
\item
  Let A be the finite field \(\mathbf{F}_{q^{2}}\) , J the involutive automorphism \(\xi \to \xi^{q}\) of A, whose \(\mathbf{F}_{q}\) is the field of invariants. If E is a vector space of dimension \(n\) over A, \(\Phi\) a non-degenerate Hermitian sesquilinear form (with respect to J) on E, show that the order of the unitary group \(\mathbf{U}(\Phi)\) is equal to
\end{enumerate}

\[
(q ^ {n} - (- 1) ^ {n}) q ^ {n - 1} (q _ {-} ^ {n - 1} - (- 1) ^ {n - 1}) q ^ {n - 2} \dots (q ^ {2} - 1) q (q + 1)
\]

(using a method analogous to that of exerc. 10 of § 5, by means of exerc. 3).

\begin{enumerate}
\def\labelenumi{\arabic{enumi})}
\setcounter{enumi}{12}
\tightlist
\item
  Let A be the finite field \(F_{q}\) (q not a multiple of 2), E a vector space of dimension n over A, and Q a non-degenerate quadratic form on E. Show that:
\end{enumerate}

\begin{enumerate}
\def\labelenumi{\alph{enumi})}
\tightlist
\item
  If n is odd, the order of the group SO(Q) is
\end{enumerate}

\[
(q ^ {n - 1} - 1) q ^ {n - 2} (q ^ {n - 3} - 1) q ^ {n - 4} \dots (q ^ {2} - 1) q.
\]

\begin{enumerate}
\def\labelenumi{\alph{enumi})}
\setcounter{enumi}{1}
\tightlist
\item
  If n = 2m is even, the order of the group SO(Q) is equal to
\end{enumerate}

\[
(q ^ {2 m - 1} - \varepsilon q ^ {m - 1}) (q ^ {2 m - 2} - 1) q ^ {2 m - 3} \dots (q ^ {2} - 1) q
\]

where \(\varepsilon = 1\) if \((-1)^{m}\Delta\) is a square in A, \(\varepsilon = -1\) otherwise, \(\Delta\) denoting the discriminant of Q with respect to an arbitrary basis of E. (Method analogous to that of exerc. 12, using exerc. 3 of §6 and exerc. 5 of chap.~V, §11.)

\begin{enumerate}
\def\labelenumi{\arabic{enumi})}
\setcounter{enumi}{13}
\tightlist
\item
  Assume that A is a commutative field, E a finite-dimensional vector space \(n \geqslant 2\) over A, \(\Phi\) a non-degenerate Hermitian sesquilinear form on E, satisfying condition (T) ( \(\S 4\) , no. 2). Show that the only endomorphisms \(w\) of E commuting with all automorphisms \(u\) belonging to the special unitary group \(\mathbf{SU}(\Phi)\) are the homotheties, except when one simultaneously has \(n = 2\) , J = 1, A being of characteristic \(\neq 2\) . (If \(n \geqslant 3\) , write that \(w\) commutes with the involutions \(u \in \mathbf{SU}(\Phi)\) , and use exerc. 3 of \(\S 4\) ; if \(n = 2\) and J \(\neq 1\) , write that \(w\) commutes with the elements of \(\mathbf{SU}(\Phi)\) whose matrix is of the form \(\begin{pmatrix} \lambda & 0 \\ 0 & \lambda^{-1} \end{pmatrix}\) with respect to an orthogonal basis of E.
\end{enumerate}

¶ 15) Let A be a commutative field of characteristic ≠ 2, E a vector space of dimension n ≥ 1 over A, Q a non-degenerate quadratic form on E. For every automorphism u ∈ O(Q), let w = u − 1, and let r be the rank of w, and W = ̄w(0).

\begin{enumerate}
\def\labelenumi{\alph{enumi})}
\item
  Show that \(w(E)\) is the subspace \(W^{0}\) orthogonal to \(W\) .
\item
  Show that if \(n=2\) , \(r=2\) , \(u\) is the product of two reflections with respect to lines of E. (Establish that if \(w(x)\) is isotropic for every non-isotropic vector \(x\in\mathrm{E}\) , \(w(x)\) is isotropic for every \(x\in\mathrm{E}\) ; we will consider separately the case where A has at least 5 elements and the case \(\mathrm{A}=\mathbf{F}_{3}\) .)
\item
  We assume n and r are arbitrary. Show that if \(w(\mathrm{E})\) is not totally isotropic, u is a product of r symmetries with respect to hyperplanes of E, and cannot be expressed as a product of a smaller number of symmetries. (Reduce to the case where W is totally isotropic, and proceed by induction on n and r. If \(W \neq \{0\}\) show that there exists a vector \(a \in W^{0}\) such that \(w(a)\) is not isotropic, arguing by contradiction and using the fact that a plane all of whose lines except at most one are isotropic is necessarily totally isotropic; then take the symmetry s with respect to the hyperplane orthogonal to \(w(a)\) , and consider the automorphism su. If \(W = \{0\}\) take \(a \in E\) such that \(w(a)\) is not isotropic, and, with the same meaning for s, consider again the automorphism su, and use b).)
\item
  We assume that \(w(E)\) is totally isotropic. If \(s\) is a symmetry with respect to a non-isotropic hyperplane H; show that the subspace of vectors invariant under \(su\) is H ∩ W, hence is of dimension \(n-r-1\) , and deduce that \(su\) cannot be the product of fewer than \(r+1\) symmetries with respect to hyperplanes. Then deduce from \(c\) ) that \(u\) is a product of \(r+2\) symmetries with respect to hyperplanes, but cannot be the product of a smaller number of symmetries.
\item
  Deduce from c) and d) that every orthogonal automorphism is a product of at most n symmetries with respect to hyperplanes.
\item
  Show that if n is odd (resp. even), for every automorphism \(u \in \mathbf{0}(\mathrm{Q})\) of determinant 1 (resp. -1), there exists a vector \(x \neq 0\) invariant under u (use e)).
\end{enumerate}

\begin{enumerate}
\def\labelenumi{\arabic{enumi})}
\setcounter{enumi}{15}
\tightlist
\item
  Given the same hypotheses as in exerc. 15, show that if \(n \geqslant 3\) , the group SO(Q) is generated by the symmetries with respect to the non-isotropic subspaces of E of dimension \(n - 2\) (reason as in prop. 5 of no. 4).
\end{enumerate}

¶ 17) The hypotheses are the same as in exerc. 15.

\begin{enumerate}
\def\labelenumi{\alph{enumi})}
\item
  Show that, for \(n \geqslant 2\) the commutator subgroup \(\Omega(Q)\) of the orthogonal group \(\mathbf{O}(Q)\) is generated by the elements \((st)^{2}\) , where \(s\) and \(t\) range over the set of symmetries with respect to hyperplanes (use prop. 5 of no. 4, and note that for any group \(\Gamma\) the subgroup generated by the squares of the elements of \(\Gamma\) contains the commutator subgroup of \(\Gamma\) ).
\item
  Show that if \(n \geqslant 3\) the commutator subgroup of \(\mathbf{SO}(\mathbf{Q})\) is generated by the squares of the elements of \(\mathbf{SO}(\mathbf{Q})\) (use exerc. 16); deduce that this group is identical to \(\Omega(\mathbf{Q})\) , and that the quotient group \(\mathbf{SO}(\mathbf{Q})/\Omega(\mathbf{Q})\) is a commutative group all of whose elements have order 2.
\item
  A plane P ⊂ E is said to be hyperbolic if it is non-isotropic and contains isotropic lines (necessarily two in number). An automorphism \(u \in \mathbf{0}(Q)\) is said to be hyperbolic if there exists a hyperbolic plane P such that \(u(x) = x\) for all \(x \in \mathrm{P}^0\) ; one then says that \(u\) is a hyperbolic transformation associated with P. Show that if Q has index \(\geqslant 1\) , every \(u \in \mathbf{0}(Q)\) is a product of hyperbolic transformations (use prop. 5 of no. 4 and exerc. 4 a) of § 4). Deduce that if P is a hyperbolic plane, every \(u \in \mathbf{0}(Q)\) can be written \(u = tv\) , where \(t\) is a hyperbolic transformation associated with P and \(v \in \Omega(Q)\) .
\end{enumerate}

¶ 18) Let A be a commutative field, E a vector space of dimension n over A, Φ a nondegenerate Hermitian sesquilinear form on E, satisfying condition (T) (§ 4, no. 2). Let V be a vector subspace of E, H \(_{v}\) the subgroup of the unitary group U(Φ) consisting of unitary automorphisms u such that u(V) = V.

\begin{enumerate}
\def\labelenumi{\alph{enumi})}
\item
  Show that, when V is not a totally isotropic subspace of dimension \(n/2\) , the image of \(\mathrm{H}_{\mathrm{v}}\) by the map \(u \to \det u\) is the subgroup of \(\mathbf{A}^*\) consisting of \(\rho \in \mathbf{A}\) such that \(\rho \bar{\rho} = 1\) .
\item
  If \(n\) is even and \(V\) is a totally isotropic subspace of dimension \(n/2\), show that the image of \(H_{v}\) by the map \(u \rightarrow \det u\) is the subgroup of A* consisting of elements of the form \(\bar{\lambda}/\lambda\) (using prop. 2 of §4, no. 2).
\item
  Let V, W be two vector subspaces of E such that the restrictions of \(\Phi\) to V and to W are equivalent. Show that there exists \(u \in \mathbf{S}\mathbf{U}(\Phi)\) such that \(u(\mathrm{V}) = \mathrm{W}\) in the following cases:
\end{enumerate}

\(^{1}\) J is distinct from the identity (use th. 3 of chap.~V, § 11, no. 5).

\(2^{\circ} \text{ J} = 1, \text{ A is of characteristic } \neq 2, \text{ V and W are not totally isotropic subspaces of dimension } n/2.\)

\begin{enumerate}
\def\labelenumi{\alph{enumi})}
\setcounter{enumi}{3}
\tightlist
\item
  Assume J = 1, A has characteristic ≠ 2, n = 2m is even, and Φ is a nondegenerate symmetric bilinear form of index m. Let V, W be two totally isotropic subspaces of dimension m in E; show that if dim(V ∩ W) = q, then for every orthogonal automorphism u such that u(V) = W, one has det u = \((-1)^{m-q}\) (using b) and prop. 2 of § 4, no. 2). Deduce that the set of totally isotropic subspaces of dimension m is the union of two intransitivity classes N \(_{1}\) , N \(_{2}\) for the group SU(Φ); if V and W are in the same class (resp. in different classes), the dimension of V ∩ W has the same parity as m (resp. does not have the same parity as m). For a similarity u (for Φ) to be direct, it is necessary and sufficient that \(u(N_{1}) = N_{1}\) (using exerc. 4 c) of § 4.)
\end{enumerate}

\begin{enumerate}
\def\labelenumi{\arabic{enumi})}
\setcounter{enumi}{18}
\tightlist
\item
  Let A be a field, E a vector space of finite dimension \textgreater{} zero over A, \(\Phi\) a non-degenerate and non-alternating \(\varepsilon\)-hermitian sesquilinear form on E. Let u be an automorphism of E such that we have
\end{enumerate}

\[
\Phi (u (x), u (x)) = \alpha \Phi (x, x)
\]

for all \(x \in \mathbf{E}\) , with \(\alpha \in \mathbf{A}\) . Show that \(u\) is a similarity with multiplier \(\alpha\) unless the following conditions hold simultaneously: A is commutative and of characteristic 2, and J is the identity (use exerc. 8 of § 1).

\begin{enumerate}
\def\labelenumi{\arabic{enumi})}
\setcounter{enumi}{19}
\item
  Let A be a field, and let L be a finite-dimensional Hermitian space over A; we assume that the metric form \(\Phi\) of L satisfies condition (T) ( \(\S\) 4, no. 2). Show that if the index of \(\Phi\) is \(>0\) there may be similarities of L, with multiplier \(\neq 1\) , and which admit no fixed point (use the reasoning of prop. 6 of no. 6, and prop. 2 of \(\S\) 4, no. 2).
\item
  Let A be a commutative field of characteristic ≠ 2, L a finite-dimensional Euclidean space over A, and Φ the metric form of L.
\end{enumerate}

\begin{enumerate}
\def\labelenumi{\alph{enumi})}
\tightlist
\item
  Show that every bijection u from L onto itself, such that
\end{enumerate}

\[
\Phi (u (x) - u (y), u (x) - u (y)) = \Phi (x - y, x - y)
\]

for all \(x, y\) in \(L\) is a displacement (use exerc. 7 of § 1).

\begin{enumerate}
\def\labelenumi{\alph{enumi})}
\setcounter{enumi}{1}
\tightlist
\item
  Show that the group of displacements is generated by the symmetries with respect to the non-isotropic hyperplanes of the affine space L (using prop. 5 of no. 4, reduce to proving that every non-isotropic translation is the product of two such symmetries).
\end{enumerate}

\begin{enumerate}
\def\labelenumi{\arabic{enumi})}
\setcounter{enumi}{21}
\item
  In a Hermitian space L, two linear varieties are said to be perpendicular if their directions are weakly orthogonal subspaces (§ 3, exerc. 11). Suppose L is of finite dimension; let \(V_{1}\) , \(V_{2}\) be two linear varieties, \(W_{1}\) , \(W_{2}\) their respective directions. Suppose that \(p = \dim(W_{1} + W_{2}) < n\) ; show that if \(W_{1} + W_{2}\) is not isotropic, there exists at least one linear variety U of dimension n - p, perpendicular to \(V_{1}\) and to \(V_{2}\) , and meeting each of the varieties \(V_{1}\) , \(V_{2}\) in exactly one point; moreover, if \(q = \dim(W_{1} \cap W_{2})\) , the union of all linear varieties U having the preceding properties is a linear variety of dimension \(n - p + q\) .
\item
  Let A be a commutative field of characteristic \(\neq 2\) , E a finite-dimensional vector space \(n + 1 \geqslant 2\) on A, Q a quadratic form on E, \(\Phi\) the symmetric bilinear form associated with Q. The set C of \(x \in E\) such that \(Q(x) = 0\) is called the isotropic cone with vertex 0 and equation \(Q(x) = 0\) . If it is not reduced to 0, the image S of \(C - \{0\}\) in the projective space \(\mathbf{P}(\mathrm{E})\) , under the canonical map \(\pi\) of \(\mathrm{E}-\{0\}\) on \(\mathbf{P}(\mathrm{E})\) (chap.~II, \(2^{\mathrm{e}}\) ed., App. III), is called a projective quadric (resp. projective conic if \(n=2\) ) with homogeneous equation \(\mathrm{Q}(x)=0\) . S is said to be degenerate if Q is degenerate. Two projective linear varieties \(\mathrm{V}_{1}, \mathrm{V}_{2}\) of \(\mathbf{P}(\mathrm{E})\) are said to be conjugate with respect to S if \(\overline{\pi}(\mathrm{V}_{1})\) and \(\overline{\pi}(\mathrm{V}_{2})\) are orthogonal (for \(\Phi\) ). The polar \(\mathrm{V}^{0}\) of a projective linear variety \(\mathrm{V}\subset\mathbf{P}(\mathrm{E})\) with respect to S is the variety such that \(\overline{\pi}(\mathrm{V}^{0})\cup\{0\}\) is the totally orthogonal subspace (for \(\Phi\) ) to \(\overline{\pi}(\mathrm{V})\cup\{0\}\) ; if V is a hyperplane and if S is non-degenerate, \(\mathrm{V}^{0}\) is reduced to a point, called the pole of V. A projective linear variety V is said to be tangent to S if \(\overline{\pi}(\mathrm{V})\cup\{0\}\) is an isotropic subspace (for \(\Phi\) ).
\end{enumerate}

We assume in what follows that S is nonempty and nondegenerate.

\begin{enumerate}
\def\labelenumi{\alph{enumi})}
\item
  Show that the intersection of S and a projective linear variety V is empty or is a quadric in V; for this quadric to be degenerate, it is necessary and sufficient that V be tangent to S.
\item
  Show that the tangent hyperplane to S at a point \(z \in S\) is the union of the lines passing through z and tangent to S.
\end{enumerate}

\[
z ^ {\prime}
\]

z with respect to a and b, that is, the point of D such that \(\begin{bmatrix} a & b \\ z' & z \end{bmatrix} = -1\) (chap.~II, 2nd ed., App. III, exerc. 4); show that \(z'\) belongs to the polar hyperplane of z with respect to S, and that there exist n of these points forming a projectively free family in P(E) and belonging to S (cf.~§ 4, exerc. 4 a)).

\begin{enumerate}
\def\labelenumi{\alph{enumi})}
\setcounter{enumi}{3}
\item
  We assume that n = 3 and that \(\Phi\) has maximum index v = 2. The set of lines contained in S is then the union of two sets \(N_{1}\) , \(N_{2}\) such that every line of \(N_{1}\) meets every line of \(N_{2}\) , but two distinct lines of \(N_{1}\) (resp. \(N_{2}\) ) do not meet (exerc. 18 d)). Let D, \(D'\) be two distinct lines belonging to \(N_{1}\) ; for every \(z \in D\) there exists one and only one line \(\Delta \in N_{2}\) passing through z; if \(u(z)\) is the point where \(\Delta\) meets \(D'\) , show that u is a projective linear map from D onto \(D'\) .
\item
  Assuming still n = 3, let D, D', D'' be three lines of P(E), no two of which intersect. Show that the union of the lines meeting D, D', and D'' is a nondegenerate quadric.
\end{enumerate}

\begin{enumerate}
\def\labelenumi{\arabic{enumi})}
\setcounter{enumi}{23}
\tightlist
\item
  The hypotheses and notation are those of exerc. 23, the quadric S being assumed nonempty and nondegenerate.
\end{enumerate}

\begin{enumerate}
\def\labelenumi{\alph{enumi})}
\item
  Show that the subgroup \(\Gamma\) of the projective group PGL(E) consisting of the projective linear bijections mapping S to itself, is the canonical image of the group of similarities with respect to Q. (Use exerc. 23 c) above, exerc. 2 a) of § 4, and exerc. 8 of § 1.)
\item
  Let \(a\) be a point of \(\mathbf{P}(\mathrm{E})\) not belonging to S, and let \(\Phi_{1}\) be the restriction of \(\Phi\) to the hyperplane orthogonal to \(\overline{\pi}(a)\) in E. Show that the subgroup of \(\Gamma\) leaving \(a\) invariant is isomorphic to the quotient of the orthogonal group \(\mathbf{U}(\Phi_1)\) by its center.
\item
  Let \(b\) be a point of S, F the (isotropic) hyperplane orthogonal to \(\overline{\pi}(b)\) in E, M a (non-isotropic) supplement of \(\overline{\pi}(b)\) with respect to F, and \(\Phi_{2}\) the restriction of \(\Phi\) to M. Show that the subgroup of \(\Gamma\) leaving \(b\) invariant is isomorphic to the group of similarities of a Euclidean space L of dimension \(n-1\) , having as metric form the inverse form ( \(\S\) 1, no. 7) of \(\Phi_{2}\) . (Note that if a similarity transformation for \(\Phi\) maps the line \(\overline{\pi}(b)\) in itself, it maps F into itself, and is entirely determined by its restriction to F).
\end{enumerate}

\begin{enumerate}
\def\labelenumi{\arabic{enumi})}
\setcounter{enumi}{24}
\tightlist
\item
  Let A be a commutative field of characteristic \(\neq 2\) , L a finite-dimensional affine space of dimension \(n \geqslant 2\) over A. We identify L with the complement of a projective hyperplane \(\mathrm{H}_{0}\) (“hyperplane at infinity”) of a projective space \(\mathbf{P}(\mathrm{E})\) of dimension \(n\) (chap.~II, \(2^{\mathrm{e}}\) ed., App. III, no. 4). A nonempty set S ⊂ L is said to be an affine quadric (resp. an affine conic if \(n = 2\) ) if S is the intersection of L and a projective quadric (resp. conic) in \(\mathbf{P}(\mathrm{E})\) (exerc. 23).
\end{enumerate}

\begin{enumerate}
\def\labelenumi{\alph{enumi})}
\item
  Show that if there exists a nondegenerate projective quadric \(\overline{\mathrm{S}}\subset\mathbf{P}(\mathrm{E})\) such that \(S=L\cap\overline{S}\) , this quadric is the only one having these properties, except when \(n=2\) , \(A=F_{3}\) and that S is reduced to 2 elements (note that apart from this exceptional case, for every point \(z\in H_{0}\) not belonging to \(\overline{S}\) there exists a line passing through \(z\) and meeting S in two distinct points). One then says that S is a nondegenerate affine quadric. Two affine linear varieties \(V_{1},V_{2}\) contained in L are said to be conjugate with respect to S if the projective linear varieties \(\overline{V}_{1},\overline{V}_{2}\) such that \(V_{i}=L\cap\overline{V}_{i}\) ( \(i=1,2\) ) are conjugate with respect to \(\overline{S}\) ; one defines similarly the polar (when it is not contained in \(H_{0}\) ) or the pole of an affine linear variety with respect to S, and the affine linear varieties tangent to S.
\item
  Assume that S is nondegenerate; show that one can choose an origin a in L such that, by identifying L in this way with a vector
\end{enumerate}

space, there exists a basis \((e_i)\) of L such that S is the set of \(x = \sum_{i=1}^{\infty} \xi_i e_i\) satisfying an equation of one of the two forms

\[
\begin{array}{c} \alpha_ {1} \xi_ {1} ^ {2} + \dots + \alpha_ {n} \xi_ {n} ^ {2} = 1 \\ \alpha_ {1} \xi_ {1} ^ {2} + \dots + \alpha_ {n - 1} \xi_ {n - 1} ^ {2} + \xi_ {n} = 0. \end{array}
\]

In the first case, the point \(a\) is well determined and is the pole with respect to \(\bar{S}\) of the hyperplane at infinity \(H_{0}\) (called the center of S). (Distinguish two cases according as \(H_{0}\) is or is not tangent to \(\bar{S}\) ; using th. 1 of § 6, no. 1 and prop. 2 of § 4, no. 2.)

\begin{enumerate}
\def\labelenumi{\arabic{enumi})}
\setcounter{enumi}{25}
\item
  Let A be an algebraically closed commutative field of characteristic \(\neq 2\) , E a finite-dimensional vector space over A, Q a non-degenerate quadratic form on E. Let \(u \in \mathbf{0}(Q)\) ; with the notations of exerc. 12 of § 4, we have G(p, p) = \{0\} except for \(p(X) = X - 1\) and \(p(X) = X + 1\) . Let M be a minimal element of the set of non-isotropic subspaces contained in \(G(p, p)\) and stable under \(u\) , and let \(p^h\) be the minimal polynomial of the restriction of \(u\) To M. Show that if \(h\) is odd, M is an indecomposable submodule of \(E_u\) and that if \(h\) is even, M is the direct sum of two isomorphic indecomposable submodules of \(E_u\) (To see that if \(h = 2k\) is even, M cannot be indecomposable, show that \(N = p^k(u)(M)\) would then be totally isotropic; if \((e_i)_{1 \leqslant i \leqslant 2k}\) is a basis of M such that \(u(e_i) = \varepsilon e_i + e_{i+1}\) for \(i \leqslant 2k - 1\) , \(u(e_{2k}) = \varepsilon e_{2k}\) (with \(\varepsilon = \pm 1\) ), show that \(e_k\) cannot be orthogonal to \(e_{k+1}\) , and deduce that the relation \(Q(u(e_k)) = Q(e_k)\) leads to a contradiction).
\item
  Let A be a commutative field of characteristic 2, E a vector space over A, of finite dimension n, Q a quadratic form on E, Φ the associated bilinear form, which is alternating, hence of even rank 2m (§ 5, no. 1, cor. 1 of th. 1).
\end{enumerate}

\begin{enumerate}
\def\labelenumi{\alph{enumi})}
\item
  Show that if \(E^{0}\) is the subspace of E (of dimension n - 2m) orthogonal to E for \(\Phi\) , one has \(Q(\lambda x + \mu y) = \lambda^{2}Q(x) + \mu^{2}Q(y)\) for all x, y in \(E^{0}\) ; in other words, the restriction \(Q_{0}\) of Q to \(E^{0}\) is a semilinear map from \(E^{0}\) (considered as a vector space over A) into A (considered as a vector space over the subfield \(A^{2}\) ), relative to the isomorphism \(\xi \to \xi^{2}\) from A to \(A^{2}\) . Let q be the dimension (over A) of the kernel \(E^{0} \cap \overline{Q}(0)\) of Q, and let \(E_{1}\) be a supplement of \(E^{0} \cap \overline{Q}(0)\) with respect to \(E^{0}\) ; we have \(n - 2m - q \leqslant [A : A^{2}]\) .
\item
  Deduce from a) that there exists a basis \((e_i)_{1 \leqslant i \leqslant n}\) of E, whose \(2m\) first vectors form a basis of a complement \(\mathbf{E}_2\) of \(\mathbf{E}^0\) in
\end{enumerate}

E, the \(n - 2m - q\) following ones a basis of \(\mathbf{E}_1\) , such that one has, for \(x = \sum_{i=1}^{n} \xi_i e_i\)

\[
\mathrm{Q} (x) = \sum_ {i = 1} ^ {m} \left(\alpha_ {i} \xi_ {i} ^ {2} + \xi_ {i} \xi_ {m + i} + \beta_ {i} \xi_ {m + i} ^ {2}\right) + \sum_ {i = 2 m + 1} ^ {n - q} \gamma_ {i} \xi_ {i} ^ {2}
\]

the \(\gamma_{i}\) ( \(2m + 1 \leqslant i \leqslant n - q\) ) being elements of A linearly independent with respect to \(A^{2}\) .

\begin{enumerate}
\def\labelenumi{\alph{enumi})}
\setcounter{enumi}{2}
\item
  We call the index of Q the maximum dimension of totally singular subspaces V of E such that \(V \cap E^{0} = \{0\}\) . Show that if v is the index of Q, one can take the basis \((e_{i})\) of E having the properties stated in b) so that \(\alpha_{i} = \beta_{i} = 0\) for \(1 \leqslant i \leqslant v\) and that the restriction of Q to the subspace of \(E_{2}\) generated by \(e_{v+1}, \ldots, e_{m}, e_{m+v+1}, \ldots, e_{2m}\) is a nondegenerate quadratic form of index 0.
\item
  We assume \(q=0\) ; let \(\mathbf{O}(\mathrm{Q})\) the group of automorphisms of E leaving Q invariant. If \(u\in\mathbf{O}(\mathrm{Q})\) , show that \(u(x)=x\) for all \(x\in\mathrm{E}^{0}\) . For every \(x\in\mathrm{E}_{2}\) , let \(u(x)=u_{0}(x)+u_{2}(x)\) , where \(u_{0}(x)\in\mathrm{E}^{0}\) and \(u_{2}(x)\in\mathrm{E}_{2}\) ; show that \(u_{2}\) belongs to the symplectic group \(\mathbf{Sp}(\Phi_{2})\) (where \(\Phi_{2}\) is the restriction of \(\Phi\) to \(\mathrm{E}_{2}\) ) and one has \(\mathrm{Q}(u_{2}(x))+\mathrm{Q}(x)\in\mathrm{Q}(\mathrm{E}^{0})\) . Conversely, for every automorphism \(u_{2}\in\mathbf{Sp}(\Phi_{2})\) such that \(\mathrm{Q}(u_{2}(x))+\mathrm{Q}(x)\in\mathrm{Q}(\mathrm{E}^{0})\) for all \(x\in\mathrm{E}_{2}\) , show that there exists one and only one linear map \(u_{0}\) of \(\mathrm{E}_{2}\) in E \(^{0}\) such that the linear map equal to \(u_{0} + u_{2}\) in E \(_{2}\) , and to the identity in E \(^{0}\) , belongs to O(Q).
\item
  Assume that A is a perfect field (A² = A) and that q = 0. Deduce from b) that every vector subspace of E of dimension ≥ 3 contains at least one vector x such that Q(x) = 0. If n is odd, one necessarily has m = v and n = 2m + 1, so that there exists a basis (ei) of E with respect to which one has
\end{enumerate}

\[
\mathrm{Q} \left(\sum_ {i = 1} ^ {n} \xi_ {i} e _ {i}\right) = \xi_ {1} \xi_ {m + 1} + \dots + \xi_ {m} \xi_ {2 m} + \xi_ {2 m + 1} ^ {2},
\]

and (with the notations of \(d\) ) \(\mathbf{O}(\mathrm{Q})\) is isomorphic to \(\mathbf{Sp}(\Phi_{2})\) ; all quadratic forms such that \(q=0\) are then equivalent. If \(n\) is even, one necessarily has \(n=2m\) , \(\nu=m\) or \(\nu=m-1\) , and there exists a basis \((e_{i})\) of E with respect to which one has

\[
\mathrm{Q} \left(\sum_ {i = 1} ^ {n} \xi_ {i} e _ {i}\right) = \xi_ {1} \xi_ {m + 1} + \dots + \xi_ {m - 1} \xi_ {2 m - 1} + \xi_ {m} \xi_ {2 m} + \lambda \left(\xi_ {m} ^ {2} + \xi_ {2 m} ^ {2}\right)\tag{1}
\]

where \(\lambda\in A\) . Let \(A_{1}\) be the field obtained by adjoining to A the roots of the polynomial \(\lambda X^{2}+X+\lambda\) show that this field is independent of the basis \((e_{i})\) with respect to which Q can be written in the form (1), and that for two quadratic forms (such that \(q=0\) ) to be equivalent, it is necessary and sufficient that the quadratic extensions of A which correspond to them in this way be identical (use Witt's theorem). Case where A is a finite field of characteristic 2.

¶ 28) Let A be a commutative field of characteristic 2, distinct from \(\mathbf{F}_{2}\) , E a vector space of dimension \(n = 2m\) on A, Q a nondegenerate quadratic form on E.

\begin{enumerate}
\def\labelenumi{\alph{enumi})}
\item
  Show that the orthogonal group \(\mathbf{O}(\mathrm{Q})\) is generated by the symmetries (which here are none other than the transvections belonging to \(\mathbf{O}(\mathrm{Q})\) ( \(\S 4\) , exerc. 6)) (reason as in exerc. 11 of \(\S 5\) ). Deduce from this that the commutator subgroup of \(\mathbf{O}(\mathrm{Q})\) is generated by the squares of the elements of \(\mathbf{O}(\mathrm{Q})\) (cf.~exerc. 17).
\item
  We assume that Q is of maximum index; let V, W be two totally singular subspaces of E (§ 4, no. 1) of dimension m. Let u be a symmetry \(x \to x + \frac{\Phi(x, a)}{\mathrm{Q}(a)} a\) (§ 4, exerc. 6); let k be the dimension of V ∩ W. Show that the dimension of V ∩ u(W) is \(k + 1\) if a is orthogonal to V ∩ W, \(k - 1\) in the contrary case (in the first case, note that \(a = x + y\) , where \(x \in V\) , \(y \in W\) and show that \(u(y) = x\) ; in the second, note that u cannot leave invariant any singular vector not orthogonal to a).
\item
  We again assume that the index of Q is arbitrary. Show that the subgroup \(\mathbf{SO}(\mathrm{Q})\) of \(\mathbf{O}(\mathrm{Q})\) consisting of the automorphisms of E that are products of an even number of symmetries, is a normal subgroup of index 2 of \(\mathbf{O}(\mathrm{Q})\) . (Show that the product of an odd number of symmetries cannot be the identity, by considering the extension of Q to the
\end{enumerate}

vector space E' obtained by extending the scalar field of E to its algebraic closure; then use b).) (Cf. § 9, exerc. 9.)

\begin{enumerate}
\def\labelenumi{\alph{enumi})}
\setcounter{enumi}{3}
\item
  If \(V_{1}\) , \(\bar{V}_{2}\) are two totally singular subspaces of E, of the same dimension \textless m, show that there exists an automorphism \(u \in \mathbf{SO}(Q)\) such that \(u(V_{1}) = V_{2}\) . Conversely, if \(V_{1}\) and \(V_{2}\) are two totally singular subspaces of dimension m, then for there to exist an automorphism \(u \in \mathbf{SO}(Q)\) such that \(u(V_{1}) = V_{2}\) , it is necessary and sufficient that the dimension of \(V_{1} \cap V_{2}\) have the same parity as m (argue as in exerc. 18, using b)).
\item
  A plane P ⊂ E is said to be hyperbolic if it is non-isotropic and contains singular lines (necessarily two in number). A transformation \(u \in \mathbf{0}(Q)\) is said to be hyperbolic if there exists a hyperbolic plane P such that \(u(x) = x\) for all \(x \in \mathrm{P}^0\) ; one then says that \(u\) is a hyperbolic transformation associated with P. Show that if Q has index \textgreater0, every \(u \in \mathbf{0}(Q)\) is a product of hyperbolic transformations (use \(a\) ). Deduce that if P is a hyperbolic plane, every transformation \(u \in \mathbf{0}(Q)\) can be written \(u = sv\) , where \(s\) is a hyperbolic transformation associated with P and \(v\) belongs to the commutator subgroup of \(\mathbf{0}(Q)\) .
\end{enumerate}

\begin{enumerate}
\def\labelenumi{\arabic{enumi})}
\setcounter{enumi}{28}
\tightlist
\item
  Given the hypotheses of exerc. 28, suppose further that Q has maximum index m; let \((e_{i})\) be a symplectic basis of E (for the alternating form \(\Phi\) associated with Q) consisting of singular vectors (§ 4, no. 2, prop. 2), so that the matrix of \(\Phi\) with respect to this basis is the matrix denoted by R in exerc. 14 of § 5. Using the notation of that exercise, show that, for a symplectic matrix \((^{t}D + S)^{-1}(D + S)\) to be the matrix of an automorphism \(u \in \mathbf{O}(\mathbb{Q})\) , it is necessary and sufficient that S be alternating (write that every vector \(u(e_{i})\) is singular, noting that one has \((^{t}D + S).u(e_{i}) = (D + S).e_{i}\) ).
\end{enumerate}

\section{Hermitian forms and ordered fields}

Throughout this paragraph, K denotes a maximal ordered field (hence commutative and of characteristic zero; cf.~Chap.~VI, § 2), and we assume that we are in one of the following three cases:

1°) A = K, J is the identity;

2°) A is the field \(\mathrm{K}(i)\) , obtained by adjoining to K a square root i of -1, and, for every \(\lambda \in A\) , \(\overline{\lambda}\) is the conjugate of \(\lambda\) (chap.~II, § 7, no. 7).

3°) A is the quaternion algebra over K corresponding to the pair \((-1,-1)\) (or, as we shall say for short, the field of quaternions over K), and for every \(\lambda\in A\) , \(\overline{\lambda}\) is the conjugate of \(\lambda\) (cf.~chap.~II, § 7, no. 8 and chap.~VIII, § 11, no. 2).

If \(\Phi\) is a Hermitian form on E, so in all cases we have, \(\Phi(x, x) \in K\) for all \(x \in E\) , since \(\Phi(x, x) = \overline{\Phi(x, x)}\) .

\subsection{Positive Hermitian forms.}

DEFINITION 1. --- A Hermitian form \(\Phi\) on E is said to be positive (resp. negative) if \(\Phi(x, x) \geqslant 0\) (resp. \(\Phi(x, x) \leqslant 0\) ) for all \(x \in \mathbf{E}\) .

When A = K, one also says that the quadratic form \(Q(x) = \frac{1}{2} \Phi(x, x)\) to which \(\Phi\) is associated is positive (resp. negative).

Suppose E is finite-dimensional over A, and let \((e_{i})\) ( \(i = 1, \ldots, n\) ) be an orthogonal basis of E ( \(\S 6, no. 1\) , th. 1). For a Hermitian form \(\Phi\) on E to be positive, it is necessary and sufficient that \(\Phi(e_{i}, e_{i}) \geqslant 0\) for \(i = 1, \ldots, n\) . Let \(\Phi\) be a nondegenerate positive Hermitian form on E; since every positive element of K is a square, hence of the form \(\rho\bar{\rho} (\rho \in A)\) , there exist in E orthonormal bases for \(\Phi\) ( \(\S 6, no. 1\) , cor. 1 of th. 1) .

PROPOSITION 1. --- Suppose E is finite-dimensional, and A = K or A = K(i). If \(\Phi\) is a nondegenerate positive Hermitian form on E, then the extensions of \(\Phi\) to \(\bigotimes^{p}\mathbf{E}\) and \(\bigwedge^{p}\mathbf{E}\) ( \(p > 0\) ), as well as the inverse form of \(\Phi\) , are nondegenerate positive Hermitian forms.

This follows immediately from the existence of an orthonormal basis of E, and from prop. 2, § 6, no. 1.

Prop. 1 remains true if one replaces everywhere “nondegenerate positive” by “positive”.

PROPOSITION 2. --- Let Φ be a positive Hermitian form on E. For x, y in E, one has

\[
\Phi (x, y) \overline {{{{\Phi (x , y)}}}} \leqslant \Phi (x, x) \Phi (y, y).\tag{1}
\]

The inequality is indeed immediate when the vectors \(x\) and \(y\) are proportional. Suppose therefore \(x\) and \(y\) linearly independent. Let A' be the (commutative) subfield K( \(\Phi(x, y)\) ) of A, F the vector plane over A' generated by \(x\) and \(y\) , and \(\Phi_{\mathbb{F}}\) the restriction of \(\Phi\) to F; this takes its values in A'. By prop. 1, the discriminant of \(\Phi_{\mathbb{F}}\) with respect to the basis \((x, y)\) is \(\geqslant 0\) . Now this discriminant is \(\Phi(x, x)\Phi(y, y) - \Phi(x, y)\overline{\Phi(x, y)}\) . QED.

COROLLARY. --- The set of isotropic vectors of E is the subspace \(E^{0}\) orthogonal to E for \(\Phi\) . For \(\Phi\) to be nondegenerate, it is necessary and sufficient that one have \(\Phi(x, x) > 0\) for all \(x \neq 0\) .

PROPOSITION 3. --- Suppose E is finite-dimensional and A is commutative (A = K or A = K(i)). Let \(\Phi\) be a Hermitian form on E, X its matrix with respect to a basis ( \(x_j\) ) ( \(j = 1, \ldots, n\) ) of E; for every subset H of {[}1, n{]}, denote by \(X_{\mathrm{H,H}}\) the minor of X obtained by deleting the rows and columns with indices \(j \notin \mathrm{H}\) (chap.~III, § 6, no. 3).

\begin{enumerate}
\def\labelenumi{\alph{enumi})}
\item
  If \(\Phi\) is nondegenerate positive, one has \(X_{\mathrm{H,H}} > 0\) for every subset H of \([1, n]\) .
\item
  Conversely, if, setting \(H_{j} = (1, j)\) , one has \(X_{H_{j}, H_{j}} > 0\) for \(j = 1, \ldots, n\) , \(\Phi\) is nondegenerate positive.
\end{enumerate}

Suppose first \(\Phi\) nondegenerate positive. The elements \((x_{j})\) ( \(j\in\mathbb{H}\) ) form a basis of a subspace F of E, and the minor \(X_{\mathrm{H,H}}\) is the discriminant of the restriction \(\Phi_{\mathrm{F}}\) of \(\Phi\) to F with respect to this basis; now, since \(\Phi_{\mathrm{F}}(x,x)>0\) for all \(x\neq0\) in F, \(\Phi_{\mathrm{F}}\) is nondegenerate positive (cor. of prop. 2); one therefore has \(X_{\mathrm{H,H}}>0\) (prop. 1). To prove \(b\) ), note that, with the notations of prop. 1, § 6, no. 1, the minor \(X_{\mathrm{Hj,Hj}}\) is equal to D \(_{j+1,j+1}\) ; there therefore exists (prop. 1, § 6, no. 1) an orthogonal basis \((e_{j})\) ( \(j=1,\ldots,n\) ) of E such that \(\Phi(e_{j},e_{j})>0\) for \(j=1,\ldots,n\) ; consequently \(\Phi\) is nondegenerate positive.

Remark. --- It follows from the existence of orthonormal bases that two nondegenerate positive Hermitian forms on two vector spaces of the same finite dimension are equivalent (§ 1, no. 6). Let then L be a finite-dimensional Hermitian space over A, whose metric form is nondegenerate positive, and V₁ and V₂ two linear varieties of the same dimension in L (§ 6, no. 6); since the restrictions of the metric form to the directions T₁ and T₂ of V₁ and V₂ on the one hand, and to the orthogonal subspaces T₁ and T₂ on the other hand, are equivalent, there exists a unitary automorphism \(u\) of the space T of translations of L such that \(u(\mathrm{T}_1) = \mathrm{T}_2\) and \(u(\mathrm{T}_1^0) = \mathrm{T}_2^0\) ; there therefore exists a displacement \(v\) of L such that \(v(\mathrm{V}_1) = \mathrm{V}_2\) . Let \((a, b), (a', b')\) two pairs of points of L; for there to exist a displacement \(v\) of L such that \(v(a) = a'\) and \(v(b) = b'\) it is necessary and sufficient (with the notation of § 6, no. 6) that one have \(e(a, b) = e(a', b')\) ; the element \(\sqrt{e(a, b)}\) of A (chap.~VI, § 2, no. 4) is called the distance of \(a\) and \(b\) in the Hermitian space L.

\subsection{The law of inertia.}

THEOREM 1 (“law of inertia”). --- Suppose that A satisfies the hypotheses at the beginning of this paragraph, and that E is of finite dimension n.~Let Φ be a Hermitian form on E. Then:

\begin{enumerate}
\def\labelenumi{\alph{enumi})}
\item
  There exists a decomposition of E as a direct sum of the subspace \(E^{0}\) orthogonal to E, and of two subspaces \(E^{+}\) and \(E^{-}\) such that the restriction of \(\Phi\) to \(E^{+}\) (resp. \(E^{-}\) ) is positive (resp. negative) and non-degenerate.
\item
  There exists an orthogonal basis \((e_i)_{1\leqslant i\leqslant n}\) of E such that
\end{enumerate}

\[
\Phi (\sum_ {i = 1} ^ {n} \xi_ {i} e _ {i}, \sum_ {i = 1} ^ {n} \eta_ {i} e _ {i}) = \sum_ {i = 1} ^ {s} \xi_ {i} \overline {{\eta}} _ {i} - \sum_ {i = s + 1} ^ {s + t} \xi_ {i} \overline {{\eta}} _ {i}.\tag{2}
\]

\begin{enumerate}
\def\labelenumi{\alph{enumi})}
\setcounter{enumi}{2}
\item
  The dimensions s of \(E^{+}\) and t of \(E^{-}\) are the same for all direct-sum decompositions satisfying the conditions stated in a); the integer s (resp. t) is the maximum of the dimensions of subspaces F of E such that the restriction of \(\Phi\) to F is positive (resp. negative) and non-degenerate.
\item
  The rank of \(\Phi\) is \(s + t\) .
\item
  If \(\Phi\) is non-degenerate, its index is equal to inf(s, t) (§ 4, no. 2, def. 2).
\end{enumerate}

Let, in fact, \((x_{i})\) ( \(i = 1, \ldots, n\) ) be an orthogonal basis of E (§ 6, no. 1, th. 1); recall that one has \(\Phi(x, x) \in K\) for all \(x \in E\) . One may assume that \(\Phi(x_{i}, x_{i}) > 0\) for \(i = 1, \ldots, s\) , \(\Phi(x_{i}, x_{i}) < 0\) for \(i = s + 1, \ldots, s + t\) and \(\Phi(x_{i}, x_{i}) = 0\) for \(i = s + t + 1, \ldots, n\) . This proves \(a\) ), for one takes as \(E^{+}\) (resp. \(E^{-}\) ) the subspace generated by \(x_{1}, \ldots, x_{s}\) (resp. by \(x_{s+1}, \ldots, x_{s+t}\) ), and \(E^{0}\) is then generated by \(x_{s+t+1}, \ldots, x_{n}\) . From this we deduce \(b\) ) by noting that \(\Phi(x_{i}, x_{i})\) is of the form \(\rho\bar{\rho}\) (resp.

\(-\rho\bar{\rho}) \quad (\rho\in\Lambda^{*}) \quad \text{for} \quad i=1,\ldots,s \quad (\text{resp. } i=s+1,\ldots,s+t).\) To prove c), consider a subspace P of E such that the restriction of \(\Phi\) to P is positive and nondegenerate; we then have \(\mathrm{P}\cap(\mathrm{E}^{-}+\mathrm{E}^{0})=\{0\}\) and the sum \(P+E^{-}+E^{0}\) is therefore direct; we conclude that dim \(P\leqslant\dim E^{+}=s\) and this proves c). Assertion d) follows immediately from a).

Finally suppose that \(\Phi\) is nondegenerate, set \(q = \inf (s, t)\) , and let us show that \(q\) is the index of \(\Phi\) . With the notations of \(b\) ), the vectors \(e_i + e_{s+i}\) (resp. \(e_i - e_{s+i}\) ) ( \(i = 1, \ldots, q\) ) generate a totally isotropic subspace F (resp. F'). Since K has characteristic 0, F + F' is generated by \(e_1, \ldots, e_q\) , \(e_{s+1}, \ldots, e_{s+q}\) and is therefore non-isotropic. The restriction of \(\Phi\) to the subspace H = (F + F')⁰ is then positive (or negative) and nondegenerate, and H therefore contains no isotropic vector \(\neq 0\) . Since F⁰ contains F and H, and dim(F + H) = codim F' = codim F, we have F⁰ = F + H. Thus an isotropic vector z orthogonal to F necessarily lies in F, because, in the direct sum F + H, the component of z in H is isotropic. Consequently, F is a maximal totally isotropic subspace, and this proves e).

DEFINITION 2. --- With the notation of Th. 1, the pair \((s,t)\) is called the signature of \(\Phi\) .

\subsection{Reduction of a form with respect to a positive Hermitian form.}

In this section we will assume that E is of finite dimension n over A, and we will denote by Φ a positive nondegenerate Hermitian form on E.

Since the linear maps from E into \(\mathbf{E}^*\) associated with \(\Phi\) are bijective, whatever the elements \(x, y, z\) of E with \(y \neq 0\) , there exists one and only one element \(t\) of E such that \(\Phi(x, z) = \overline{\Phi(t, y)}\) In particular, if \(\mathrm{A} = \mathrm{K}\) or \(\mathrm{A} = \mathrm{K}(i)\) and if \(u\) is a semilinear map for J from E into itself (chap.~II, App. I, no. 1), then for every \(x \in \mathrm{E}\) there exists one and only one element \(u^*(x)\) such that for all \(y \in \mathrm{E}\) , \(\Phi(x, u(y)) = \overline{\Phi(u^*(x), y)}\) We see immediately that \(u^*\) is a semilinear map for J from E into itself; it is called the adjoint of \(u\) .

Remarks. --- 1) When J is the identity, we recover the notion of the adjoint of a homomorphism defined in § 1, no. 8.

\begin{enumerate}
\def\labelenumi{\arabic{enumi})}
\setcounter{enumi}{1}
\tightlist
\item
  Assume always that A = K or A = K(i). Let E \(^{j}\) be the right A-module defined in § 1, no. 2, def. 5. The form Φ is a bilinear form on E × E \(^{j}\) , Φ \(^{j}\) is a bilinear form on E \(^{j}\) × E, and u is an A-linear map from E to E \(^{j}\) . The adjoint of this homomorphism, in the sense of § 1, no. 8, is an A-linear map from E into E \(^{j}\) ; we see at once that this coincides with the map u* defined above.
\end{enumerate}

We say that an endomorphism \(u\) of E is normal (for \(\Phi\) ) if one has \(uu^{*}=u^{*}u\) .
